\section{An open cone of limiting interior}\label{sec:noncollapse}

The excluded cone supplied by Proposition~\ref{prop:boundary-data} could,
at this stage, be all of $\R^k$.  We rule this out by constructing actual
interior balls whose radii are comparable to their distance from the chosen
boundary point.  The analytic interior graphs and the continuity principle
are the essential inputs; the supporting functions will enter only in the
next section.
The construction uses the analytic-disc method introduced by
Bishop~\cite{Bishop1965}. In particular, the nonlinear boundary equation
expresses the real part in terms of the imaginary part through harmonic
conjugation; see Hill--Taiani~\cite[Section~2, Equation~(2.1) and
Proposition~2.1, pp.~330--331]{HillTaiani1978} for its
classical formulation. We prove here the parameter estimates and the
regularization needed to obtain included balls from the approaching graphs.

We first isolate the analytic-disc construction.  All balls in this section
are Euclidean balls, with the real or complex ambient space indicated by
their centers.  In the application, we fix a complex-normal slice through
the chosen boundary point and straighten its totally real boundary graph.
The approaching interior graphs then satisfy the hypotheses of the lemma
below.  A ball in this slice will give a ball in a normal-offset fiber,
which the joint limits of Section~\ref{sec:preparation} can detect.

\begin{lemma}[Interior balls from approaching graphs]\label{lem:disc-family}
Let $\Omega$ be open in a neighborhood of $0\in\C^k$.  Assume that the
continuity principle of Lemma~\ref{lem:continuity} holds for continuous
deformations of closed analytic discs contained in that neighborhood.
Suppose that $f(X,u)$ is an $\R^k$-valued real-analytic function for real
$X$ near $0$ and real $u$ near $0$, that $f(X,0)=0$, and that
\[
 X+i f(X,u)\in\Omega
 \qquad\text{for all such $X$ and all sufficiently small $u>0$.}
\]
Then there are constants $c,C,t_0>0$ and points $a_t\in\C^k$, for
$0<t<t_0$, such that
\begin{equation}\label{eq:disc-balls}
 |a_t|\le Ct,
 \qquad B(a_t,ct)\subset\Omega.
\end{equation}
The discs used to obtain these balls stay in an arbitrarily small fixed
neighborhood of $0$ after the graph neighborhood is restricted.
\end{lemma}

\begin{proof}
Write $f_u(X)=f(X,u)$.  Choose nonnegative smooth functions $\lambda_0$
and $\eta$ on the unit circle, both zero on a fixed arc about $1$, with
$\eta\not\equiv0$ and
\[
 \operatorname{supp}\eta\Subset\{\lambda_0>0\}.
\]
Set $\lambda=\alpha\lambda_0$, where the positive constant $\alpha$ will
be fixed sufficiently small.  For parameters $b,e\in\R^k$ and
$\delta\ge0$, we seek a holomorphic disc $A_{b,e,\delta}$ whose boundary
values $X+iY$ satisfy
\begin{equation}\label{eq:disc-boundary}
 X(0)=b,
 \qquad
 Y(\theta)=f_{\lambda(\theta)+\delta}(X(\theta))+
             \eta(\theta)e.
\end{equation}
Here $X(0)$ means the value at the boundary angle $\theta=0$, not the
value at the center of the disc.

The parameters $b$ and $e$ have different geometric roles.  When $\delta=0$,
the boundary arc about $1$ maps into $\R^k$, with endpoint
$A_{b,e,0}(1)=b$.  Varying $b$ moves that endpoint in the real directions;
the perturbation $\eta e$, supported away from $1$, will change the inward
radial derivative in all imaginary directions.  For $B\in\R^k$,
evaluation at $1-t$ with $b=tB$ will therefore give independent
displacements of order $t$ in all
$2k$ real directions.  We first prove this derivative assertion and then
use a positive $\delta$ to put the evaluated discs inside $\Omega$.

\paragraph{Solving the boundary equation.}
Let $\mathcal H=H^2(S^1,\R^k)$ be the real Sobolev space of functions
with two square-integrable derivatives.  The Fourier-series estimate
\[
 \sum_{j\in\Z}\frac{j^2}{(1+j^2)^2}<\infty
\]
gives the continuous embedding $\mathcal H\hookrightarrow C^1(S^1)$.
This space is a Banach algebra under multiplication.  Let $T$ be the
bounded conjugation operator which recovers the real part of a
holomorphic boundary value from its imaginary part, with real part of
mean zero, and put
\[
 \mathcal T Y=TY-(TY)(0).
\]
Thus $\mathcal T$ is bounded on $\mathcal H$, and
Equation~\eqref{eq:disc-boundary} is equivalent to
\begin{equation}\label{eq:disc-contraction}
 X=b+\mathcal T\bigl(f_{\lambda+\delta}(X)+\eta e\bigr).
\end{equation}

Here is the smallness needed for contraction.  The identity $f(X,0)=0$
gives a real-analytic factorization $f(X,u)=u\widetilde f(X,u)$.  On a fixed small
$\mathcal H$-ball of functions $X$, the product rule through two
derivatives, the embedding into $C^1$, and bounded derivatives of $\widetilde f$
give
\begin{equation}\label{eq:disc-composition-bound}
 \|f_{\lambda+\delta}(X)\|_{\mathcal H}
 +\|D_X[f_{\lambda+\delta}(X)]\|_{\mathcal L(\mathcal H)}
 \le C_0(\alpha+|\delta|).
\end{equation}
For example, both expressions contain the factor $\lambda+\delta$;
derivatives falling on that factor are also $O(\alpha+|\delta|)$.
The remaining composition and multiplication operators are uniformly
bounded on the chosen ball.  For all sufficiently small $\alpha$ and
$|\delta|$, the Lipschitz constant in
Equation~\eqref{eq:disc-contraction} is less than $1/2$.  On common
sufficiently small $b,e$ neighborhoods, its right side maps the chosen
ball into itself.  The
contraction theorem gives a unique solution $X$.

Smooth composition on $H^2$, or its direct proof by the same product
estimates, shows that this equation is smooth in $X,b,e,\delta$.
Its $X$-linearization is the identity minus an operator of norm less
than $1/2$.  The implicit function theorem therefore gives smooth
parameter dependence in $\mathcal H$.  In particular $\partial_e X$
is uniformly bounded as $\alpha$ decreases; it need not be small.
The boundary function $X+iY$ extends holomorphically to the disc and
continuously to its closure.  All these conclusions hold also with
$\lambda$ replaced by $\rho\lambda$, $0\le\rho\le1$.  Shrinking the
fixed $\mathcal H$-ball keeps every disc in the prescribed coordinate
neighborhood, by the boundary bound and the maximum principle.

\paragraph{The limiting radial derivative.}
For $\delta=0$, the function $Y$ vanishes on the fixed arc about
$\theta=0$.  Hence $A_{b,e,0}(1)=b$, and reflection across that arc
shows that the radial derivative there is purely imaginary.  More
explicitly, with $Y=Y_{b,e,0}$, the Schwarz integral gives, near $\zeta=1$,
\[
 A_{b,e,0}(\zeta)-b
 =\frac{i}{2\pi}\int_{-\pi}^{\pi}
 \left(\frac{e^{i\theta}+\zeta}{e^{i\theta}-\zeta}
       -\frac{e^{i\theta}+1}{e^{i\theta}-1}\right)
 Y(\theta)\,d\theta.
\]
There is no singularity on the support of $Y$.  Expanding the kernel
at $\zeta=1$ gives
\begin{align}
 A_{b,e,0}(1-t)
   &=b+i t D(b,e)+O(t^2),\label{eq:disc-radial}\\
 D(b,e)
   &=\frac{1}{2\pi}\int_{-\pi}^{\pi}
        \frac{2Y_{b,e,0}(\theta)}{|1-e^{i\theta}|^2}\,d\theta.
        \label{eq:disc-normal-derivative}
\end{align}
The remainder is uniform with its first derivatives in $b,e$ on a
fixed smaller parameter neighborhood.  This follows from the same
integral formula and the smooth $\mathcal H$-dependence proved above.

Differentiating Equation~\eqref{eq:disc-normal-derivative} in $e$ gives
\[
 \partial_e D(b,e)=c_\eta I+R_\alpha(b,e),
 \qquad
 c_\eta=\frac{1}{2\pi}\int_{-\pi}^{\pi}
       \frac{2\eta(\theta)}{|1-e^{i\theta}|^2}\,d\theta>0.
\]
The remainder contains $\partial_X f_{\lambda(\theta)}(X)\partial_e X$.
The first factor is $O(\alpha)$ and the second is uniformly bounded.
Their product vanishes on the arc where $\lambda=0$, so the integral
has a uniformly bounded kernel on its support.  Thus
$\|R_\alpha\|\le C_1\alpha$.  The positive constant $c_\eta$ is
independent of $\alpha$.  Fix $\alpha$ sufficiently small that
$\partial_e D(0,0)$ is invertible, and only then fix the smaller
$b,e,\delta$ parameter neighborhoods ensuring all the conclusions above.

\paragraph{Interior discs for positive regularization.}
For this fixed positive $\alpha$, on $\operatorname{supp}\eta$
the function $\lambda$ has a positive minimum.  As $X$ ranges in a
fixed smaller closed real ball compactly contained in the graph chart,
$\theta$ ranges in this support, and $\delta$
ranges in a sufficiently small closed interval $[0,\delta_0]$, the
points
\[
 X+i f_{\lambda(\theta)+\delta}(X)
\]
form a compact subset of $\Omega$.  Consequently, after reducing the
allowed size of $e$, their displacements by $i\eta(\theta)e$ still
belong to $\Omega$, uniformly for $0\le\delta\le\delta_0$.
Outside this support there is no displacement.  Thus for every
$\delta>0$ all boundary values in Equation~\eqref{eq:disc-boundary}
are in $\Omega$.

For a fixed $\delta>0$, set $e=0$ and increase $\rho$ from $0$ to $1$
in $\rho\lambda$.  At $\rho=0$ the solution is the constant disc
$b+i f_\delta(b)$.  All boundary values during this deformation lie
on interior graphs, so the continuity principle puts every disc in
$\Omega$.  Next increase the displacement from $0$ to $e$, now with
$\lambda$ fixed.  The same principle applies by the preceding compact
interior bound.  We have proved
\begin{equation}\label{eq:positive-discs}
 A_{b,e,\delta}(\overline\D)\subset\Omega
 \qquad(\delta>0).
\end{equation}
No inclusion is claimed for $\delta=0$.
The final radius allowed for $e$ may depend on the fixed $\alpha$.
No uniform lower bound as $\alpha\downarrow0$ is needed: the included-ball
constants below are fixed only after these choices.

\paragraph{Interior balls.}
For real parameters $(B,e)$ in a fixed small ball in $\R^{2k}$, define
\[
 \Psi_t^0(B,e)=t^{-1}A_{tB,e,0}(1-t).
\]
Equation~\eqref{eq:disc-radial} gives convergence in $C^1$ on that ball:
\begin{equation}\label{eq:disc-scaled-limit}
 \Psi_t^0(B,e)\longrightarrow\Psi_0(B,e):=B+iD(0,e).
\end{equation}
The real differential $J=D\Psi_0(0,0)$ is invertible.
Smooth parameter dependence and bounded evaluation on the closed disc
also give, on a fixed compact parameter set,
\[
 \|A_{b,e,\delta}-A_{b,e,0}\|
       _{C^1((b,e);\,C^0(\overline\D))}\le C_2\delta.
\]
Indeed, integrate the uniformly bounded derivatives $\partial_\delta A$,
$\partial_b\partial_\delta A$, and $\partial_e\partial_\delta A$
over $[0,\delta]$; their bounds in $H^2$ control the displayed norm.
We may therefore take the explicit positive regularization
$\delta(t)=t^2$ for sufficiently small $t$.  The maps
\[
 \Psi_t(B,e)=t^{-1}A_{tB,e,t^2}(1-t)
\]
differ from $\Psi_t^0$ by $O(t)$ in $C^1$: the value and $e$-derivative
errors are $O(\delta/t)$, while the $B$-derivative error is $O(\delta)$.
They therefore satisfy the same limit
in Equation~\eqref{eq:disc-scaled-limit}.  By
Equation~\eqref{eq:positive-discs}, their images lie in $t^{-1}\Omega$.

For completeness, this $C^1$ convergence gives included balls, not
merely interior limit points.  Choose a closed parameter ball
$\overline B(0,r)$ so small that
$\|J^{-1}(D\Psi_0-J)\|<1/4$ there.  For small $t$ the corresponding
bound for $\Psi_t$ is less than $1/2$.  If
\[
 |y-\Psi_t(0)|<\frac{r}{2\|J^{-1}\|},
\]
the map $x\mapsto x-J^{-1}(\Psi_t(x)-y)$ is a contraction of
$\overline B(0,r)$ into itself.  Its fixed point proves that $y$ is
in the image of $\Psi_t$.  These images thus contain balls of one
fixed positive radius, centered at the bounded points $\Psi_t(0)$.
Multiplication by $t$ proves Equation~\eqref{eq:disc-balls}.
\end{proof}

\begin{proposition}[Noncollapse]\label{prop:noncollapse}
At the point selected in Proposition~\ref{prop:supported-point}, the
excluded cone in Proposition~\ref{prop:boundary-data} satisfies
$F\ne\R^k$.
\end{proposition}

\begin{proof}
Keep the old coordinates $(z_0,w_0)$ and fix the complex slice
$z_0=z_0(p)$.  The graph of $S$ in this slice is totally real and real
analytic.  Its parametrization by $\operatorname{Re}w_0$ therefore
has invertible complexified derivative, and its complexification
provides a local biholomorphic change of coordinates sending this
graph to $\R^k$ and $p$ to $0$.

Restrict the analytic interior graphs from
Proposition~\ref{prop:boundary-data} to this slice.  In the new
coordinates their real-part projection is the identity at $u=0$.
The real-analytic implicit function theorem consequently writes them,
on a common neighborhood, as
\[
 X+i f_u(X),\qquad f_0(X)=0.
\]
The identity at $u=0$ holds for every $X$ in this neighborhood, not
just for $X=0$.  For $u>0$ these graphs lie in the slice of $U$.
The continuity principle in the ambient smooth chart applies to
discs in this fixed slice, and is unchanged by its biholomorphic
coordinate change.  Lemma~\ref{lem:disc-family} now gives slice balls
of radii $ct$ with centers $O(t)$ from $p$.

Return to the $w_0$-coordinates.  Uniform bounds for the coordinate
change and its inverse, after one fixed shrinking, give constants
$c',C'>0$ and centers $w_t$ such that
\[
 B(w_t,c't)\subset\{w_0:(z_0(p),w_0)\in U\},
 \qquad |w_t-w_0(p)|\le C't.
\]
Set
\[
 s_t=(z_0(p),\operatorname{Re}w_t),
 \qquad v_t=\operatorname{Im}w_t-\phi(s_t).
\]
Intersecting each complex ball with the real affine slice
$\operatorname{Re}w_0=\operatorname{Re}w_t$ gives
\begin{equation}\label{eq:offset-balls}
 B(v_t,c't)\subset A_{s_t},
 \qquad s_t\longrightarrow s_*,
 \qquad |v_t|=O(t).
\end{equation}
The last estimate uses the smoothness of $\phi$ and the fact that
$p$ lies on its graph.

Choose a sequence $t_j\downarrow0$ for which $v_{t_j}/t_j\to v$.
Equation~\eqref{eq:offset-balls} implies
\[
 \operatorname{dist}\bigl(v,\R^k\setminus t_j^{-1}A_{s_{t_j}}\bigr)
 \ge c'-|v-v_{t_j}/t_j|\ge c'/2
\]
for large $j$.  If $v$ belonged to $F$, the moving-parameter
excluded-offset conclusion of Proposition~\ref{prop:boundary-data}
would force this distance to tend to zero.  Thus $v\notin F$, as
required.
\end{proof}
